% Part: incompleteness
% Chapter: representability-in-q
% Section: minimization-representable

\documentclass[../../../include/open-logic-section]{subfiles}

\begin{document}

\olfileid[id]{inc}{req}{min}
\olsection{Pencarian Tak Berbatas pada Fungsi Reguler Dapat
Direpresentasikan dalam $\Th{Q}$}

Mari kita perhatikan pencarian tak berbatas. Andaikan $g(x, z)$ bersifat
reguler dan direpresentasikan dalam $\Th{Q}$, misalnya oleh
!!{formula}~$!A_g(x, z, y)$. Misalkan $f$ didefinisikan oleh
$f(z) = \umin{x}{[g(x, z) = 0]}$. Kita ingin menemukan
!!a{formula}~$!A_f(z, y)$ yang merepresentasikan~$f$. Nilai~$f(z)$ ialah
bilangan~$x$ yang (a) memenuhi $g(x, z) = 0$ dan (b) merupakan bilangan
terkecil semacam itu, yakni untuk setiap $w < x$, $g(w, z) \neq 0$. Jadi,
berikut merupakan pilihan yang alami:
\[
!A_f(z,y) \ident !A_g(y, z, \Obj 0) \land \lforall[w][(w < y \lif \lnot
  !A_g(w, z, \Obj 0))].
\]
Dalam kasus umum, tentu saja, kita harus mengganti $z$ dengan $z_0$, \dots,
$z_k$.

Pembuktiannya kembali melibatkan beberapa lema mengenai pernyataan-pernyataan
yang dapat dibuktikan oleh $\Th{Q}$ karena teori tersebut cukup kuat.

\begin{lem}
\ollabel{lem:succ} Untuk setiap !!{constant}~$a$ dan setiap bilangan
asli~$n$,
\[
\Th{Q} \Proves \eq[(a' + \num n)][(a + \num n)'].
\]
\end{lem}

\begin{proof}
Seperti biasa, pembuktian dilakukan dengan induksi pada $n$. Dalam kasus
dasar, $n = 0$, kita perlu menunjukkan bahwa $\Th{Q}$ membuktikan
$\eq[(a' + \Obj 0)][(a + \Obj 0)']$. Namun, kita mempunyai:
\begin{align}
  \Th{Q} & \Proves \eq[(a' + \Obj 0)][a'] \quad
  \text{menurut aksioma $Q_4$} \ollabel{step1}\\
  \Th{Q} & \Proves  \eq[(a + \Obj 0)][a] \quad
  \text{menurut aksioma $Q_4$} \ollabel{step2} \\
  \Th{Q} & \Proves \eq[(a + \Obj 0)'][a'] \quad
  \text{menurut \olref{step2}} \ollabel{step3} \\
  \Th{Q} & \Proves \eq[(a' + \Obj 0)][(a + \Obj 0)'] \quad
  \text{menurut \olref{step1} dan \olref{step3}}\notag
\end{align}
Dalam langkah induksi, kita dapat mengasumsikan bahwa telah kita tunjukkan
$\Th{Q} \Proves \eq[(a' + \num n)][(a + \num n)']$. Karena
$\num{n+1}$ ialah $\num{n}'$, kita perlu menunjukkan bahwa $\Th{Q}$
membuktikan $\eq[(a' + \num{n}')][(a + \num{n}')']$. Kita mempunyai:
\begin{align}
  \Th{Q} & \Proves \eq[(a' + \num n')][(a' + \num n)'] \quad
  \text{menurut aksioma $!Q_5$} \ollabel{step5}\\
  \Th{Q} & \Proves \eq[(a' + \num n)'][(a + \num n)''] \quad
  \text{menurut hipotesis induksi} \ollabel{step6}\\
  \Th{Q} & \Proves \eq[(a + \num n')'][(a + \num n)''] \quad
  \text{menurut aksioma $!Q_5$} \ollabel{step7}\\
  \Th{Q} & \Proves \eq[(a' + \num n')][(a + \num n')'] \quad
  \text{menurut \olref{step5}, \olref{step6}, dan \olref{step7}.} \notag
\end{align}
\end{proof}

Perlu kembali disebutkan bahwa hal ini lebih lemah daripada mengatakan bahwa
$\Th{Q}$ membuktikan $\lforall[x][\lforall[y][(x' + y) = (x + y)']]$.
Walaupun !!{sentence} ini benar dalam~$\Struct{N}$, $\Th{Q}$ tidak
membuktikannya.

\begin{lem}
\ollabel{lem:less-zero}
$\Th{Q} \Proves \lforall[x][\lnot x < \Obj 0]$.
\end{lem}

\begin{proof}
Kita memberikan pembuktian secara informal (yakni hanya memberi petunjuk
tentang cara menyusun !!{derivation} formalnya).

Kita harus membuktikan $\lnot a < \Obj 0$ untuk~$a$ sebarang. Menurut
definisi $<$, kita perlu membuktikan
$\lnot \lexists[y][\eq[(y'+a)][\Obj 0]]$ dalam $\Th{Q}$. Kita akan
mengasumsikan $\lexists[y][\eq[(y'+a)][\Obj 0]]$ lalu membuktikan suatu
kontradiksi. Andaikan $\eq[(b'+a)][\Obj 0]$. Dengan menggunakan $!Q_3$, kita
mempunyai $\eq[a][\Obj 0] \lor \lexists[y][\eq[a][y']]$. Kita bedakan
kasus-kasusnya.

Kasus 1: $\eq[a][\Obj 0]$ berlaku. Dari $\eq[(b'+a)][\Obj 0]$, kita
mempunyai $\eq[(b' + \Obj 0)][\Obj 0]$. Menurut aksioma~$!Q_4$ dari
$\Th{Q}$, kita mempunyai $\eq[(b' + \Obj 0)][b']$, sehingga
$\eq[b'][\Obj 0]$. Namun, menurut aksioma~$!Q_2$, kita juga mempunyai
$\eq/[b'][\Obj 0]$, suatu kontradiksi.

Kasus 2: Untuk suatu $c$, $\eq[a][c']$. Namun, kita kemudian mempunyai
$\eq[(b' + c')][\Obj 0]$. Menurut aksioma~$!Q_5$, kita mempunyai
$\eq[(b' + c)'][\Obj 0]$, yang kembali bertentangan dengan aksioma~$Q_2$.
\end{proof}

\begin{lem}
\ollabel{lem:less-nsucc}
Untuk setiap bilangan asli~$n$,
  \[
  \Th{Q} \Proves
  \lforall[x][(x < \num {n+1} \lif (\eq[x][\Obj 0] \lor \dots \lor
    \eq[x][\num n]))].
  \]
\end{lem}

\begin{proof}
Kita menggunakan induksi pada~$n$. Mari kita perhatikan kasus dasar ketika
$n = 0$. Dalam kasus itu, kita perlu menunjukkan
$a < \num 1 \lif \eq[a][\Obj 0]$, untuk~$a$ sebarang. Andaikan
$a < \num 1$. Maka, menurut aksioma yang mendefinisikan $<$, kita mempunyai
$\lexists[y][\eq[(y'+a)][\Obj 0']]$ (karena
$\num 1 \ident \Obj 0'$).

Andaikan $b$ mempunyai sifat tersebut, yakni kita mempunyai
$\eq[(b'+a)][\Obj 0']$. Kita perlu menunjukkan $\eq[a][\Obj 0]$. Menurut
aksioma~$!Q_3$, kita mempunyai $\eq[a][\Obj 0]$ atau terdapat suatu $c$
sedemikian sehingga $\eq[a][c']$. Dalam kasus pertama, tidak ada yang perlu
ditunjukkan. Jadi, andaikan $\eq[a][c']$. Maka kita mempunyai
$\eq[(b' + c')][\Obj 0']$. Menurut aksioma~$!Q_5$ dari $\Th{Q}$, kita
mempunyai $\eq[(b'+c)'][\Obj 0']$. Menurut aksioma~$!Q_1$, kita mempunyai
$\eq[(b' + c)][\Obj 0]$. Akan tetapi, menurut aksioma~$!Q_8$, ini berarti
bahwa $c < \Obj 0$, bertentangan dengan \olref{lem:less-zero}.

Sekarang untuk langkah induksi. Kita membuktikan kasus bagi $n+1$ dengan
mengasumsikan kasus bagi~$n$. Jadi, andaikan $a < \num {n+2}$. Dengan
kembali menggunakan $!Q_3$, kita dapat membedakan dua kasus:
$\eq[a][\Obj 0]$, dan untuk suatu $b$, $\eq[a][b']$. Dalam kasus pertama,
$\eq[a][\Obj 0] \lor \dots \lor \eq[a][\num{n+1}]$ mengikuti secara
langsung. Dalam kasus kedua, kita mempunyai $b' < \num {n+2}$, yakni
$b' < \num{n+1}'$. Menurut aksioma~$!Q_8$, untuk suatu $c$,
$\eq[(c'+b')][\num{n+1}']$. Menurut aksioma $!Q_5$,
$\eq[(c'+b)'][\num{n+1}']$. Menurut aksioma~$!Q_1$,
$\eq[(c'+b)][\num{n+1}]$, sehingga $b < \num{n+1}$ menurut
aksioma~$!Q_8$. Menurut hipotesis induksi,
$\eq[b][\Obj 0] \lor \dots \lor \eq[b][\num{n}]$. Dari sini, berdasarkan
logika kita memperoleh
$\eq[b'][\Obj 0'] \lor \dots \lor \eq[b'][\num{n}']$, sehingga
$\eq[a][\num{1}] \lor \dots \lor \eq[a][\num{n+1}]$ karena
$\eq[a][b']$.
\end{proof}

\begin{lem}
  \ollabel{lem:trichotomy} Untuk setiap bilangan asli~$m$,
  \[
  \Th{Q} \Proves
  \lforall[y][((y < \num{m} \lor \num{m} < y) \lor \eq[y][\num{m}])].
  \]
\end{lem}

\begin{proof}
Dengan induksi pada~$m$. Pertama, perhatikan kasus $m=0$.
$\Th{Q} \Proves
\lforall[y][(\eq[y][\Obj 0] \lor \lexists[z][\eq[y][z']])]$ menurut~$!Q_3$.
Misalkan $a$ sebarang. Maka berlaku $\eq[a][\Obj 0]$, atau untuk suatu~$b$,
$\eq[a][b']$. Dalam kasus pertama, kita juga mempunyai
$(a < \Obj 0 \lor \Obj 0 < a) \lor \eq[a][\Obj 0]$. Namun, jika
$\eq[a][b']$, maka $\eq[(b' + \Obj 0)][(a + \Obj 0)]$ berdasarkan logika
dari~$\eq$. Menurut $!Q_4$, $\eq[(a + \Obj 0)][a]$, sehingga kita mempunyai
$\eq[(b' + \Obj 0)][a]$, dan oleh karena itu
$\lexists[z][\eq[(z' + \Obj 0)][a]]$. Menurut definisi $<$ dalam $!Q_8$,
$\Obj 0 < a$. Jika $\Obj 0 < a$, maka juga
$(\Obj 0 < a \lor a < \Obj 0) \lor \eq[a][\Obj 0]$.

Sekarang andaikan kita mempunyai
\begin{align*}
  \Th{Q} & \Proves \lforall[y][((y < \num{m} \lor \num{m} < y) \lor
    \eq[y][\num{m}])]
  \intertext{dan kita ingin menunjukkan}
  \Th{Q} & \Proves \lforall[y][((y < \num{m+1} \lor \num{m+1} < y) \lor
    \eq[y][\num{m+1}])]
\end{align*}
Misalkan $a$ sebarang. Menurut $!Q_3$, berlaku $\eq[a][\Obj 0]$ atau, untuk
suatu~$b$, $\eq[a][b']$. Dalam kasus pertama, kita mempunyai
$\eq[\num{m}' + a][\num{m+1}]$ menurut $!Q_4$, sehingga
$a < \num{m+1}$ menurut $!Q_8$.

Sekarang perhatikan kasus kedua, $\eq[a][b']$. Menurut hipotesis induksi,
$(b < \num{m} \lor \num{m} < b) \lor \eq[b][\num{m}]$.

Disjungsi pertama $b < \num{m}$ ekuivalen (menurut $!Q_8$) dengan
$\lexists[z][\eq[(z' + b)][\num{m}]]$. Andaikan $c$ mempunyai sifat ini.
Jika $\eq[(c' + b)][\num{m}]$, maka juga
$\eq[(c' + b)'][\num{m}']$. Menurut $!Q_5$,
$\eq[(c' + b)'][(c' + b')]$. Dengan demikian,
$\eq[(c' + b')][\num{m}']$. Kita memperoleh
$\lexists[u][\eq[(u' + b')][\num{m+1}]]$ dengan menggeneralisasi secara
eksistensial atas~$c$ dan mengingat bahwa
$\num{m}' \ident \num{m+1}$. Jadi, jika $b < \num{m}$, maka
$b' < \num{m+1}$ dan dengan demikian $a < \num{m+1}$.

Sekarang andaikan $\num{m} < b$, yakni
$\lexists[z][\eq[(z' + \num{m})][b]]$. Andaikan $c$ merupakan~$z$ semacam
itu, yakni $\eq[(c' + \num{m})][b]$. Berdasarkan logika,
$\eq[(c' + \num{m})'][b']$. Menurut $!Q_5$,
$\eq[(c' + \num{m}')][b']$. Karena $\eq[a][b']$ dan
$\num{m}' \ident \num{m+1}$, kita mempunyai
$\eq[(c' + \num{m+1})][a]$. Menurut $!Q_8$, $\num{m+1} < a$.

Terakhir, andaikan $\eq[b][\num{m}]$. Maka, berdasarkan logika,
$\eq[b'][\num{m}']$, sehingga $\eq[a][\num{m+1}]$.

Jadi, dari setiap disjungsi dalam kasus bagi~$m$ dan~$b$, kita dapat
memperoleh disjungsi yang bersesuaian bagi~$m+1$ dan~$a$.
\end{proof}

\begin{prop}
\ollabel{prop:rep-minimization}
Jika $!A_g(x, z, y)$ merepresentasikan $g(x, z)$ dalam~$\Th{Q}$, maka
\[
!A_f(z,y) \ident !A_g(y, z, \Obj 0) \land \lforall[w][(w < y \lif \lnot
  !A_g(w, z, \Obj 0))]
\]
merepresentasikan $f(z) = \umin{x}{[g(x, z) = 0]}$.
\end{prop}

\begin{proof}
Pertama, kita menunjukkan bahwa jika $f(n) = m$, maka
$\Th{Q} \Proves !A_f(\num n, \num m)$, yakni
\begin{align}
  \Th{Q} & \Proves !A_g(\num{m}, \num{n}, \Obj 0) \land
  \lforall[w][(w < \num{m} \lif \lnot !A_g(w, \num{n}, \Obj 0))]. \notag
  \intertext{Karena $!A_g(x, z, y)$ merepresentasikan $g(x, z)$ dan
    $g(m, n) = 0$ jika $f(n) = m$, kita mempunyai}
\Th{Q} & \Proves !A_g(\num{m}, \num{n}, \Obj 0). \notag
\intertext{Jika $f(n) = m$, maka untuk setiap $k < m$, $g(k, n) \neq 0$.
  Jadi,}
\Th{Q} & \Proves \lnot !A_g(\num{k}, \num{n}, \Obj 0). \notag
\intertext{Kita memperoleh}
\Th{Q} & \Proves \lforall[w][(w < \num{m} \lif \lnot
  !A_g(w, \num{n}, \Obj 0))]. \ollabel{rep-less}
\end{align}
menurut \olref{lem:less-zero} ketika $m = 0$, dan menurut
\olref{lem:less-nsucc} dalam kasus lainnya.

Sekarang mari kita tunjukkan bahwa jika $f(n) = m$, maka
$\Th{Q} \Proves
\lforall[y][(!A_f(\num{n}, y) \lif \eq[y][\num{m}])]$. Kita kembali
menguraikan argumennya secara informal dan menyerahkan formalisasinya kepada
pembaca.

Andaikan $!A_f(\num{n}, b)$. Dari sini kita memperoleh (a)
$!A_g(b, \num{n}, \Obj 0)$ dan (b)
$\lforall[w][(w < b \lif \lnot !A_g(w, \num{n}, \Obj 0))]$. Menurut
\olref{lem:trichotomy},
$(b < \num{m} \lor \num{m} < b) \lor \eq[b][\num{m}]$. Kita akan
menunjukkan bahwa baik $b < \num{m}$ maupun $\num{m} < b$ menghasilkan
kontradiksi.

Jika $\num{m} < b$, maka dari~(b) diperoleh
$\lnot !A_g(\num{m}, \num{n}, \Obj 0)$. Namun, $m = f(n)$, sehingga
$g(m, n) = 0$, dan dengan demikian
$\Th{Q} \Proves !A_g(\num{m}, \num{n}, \Obj 0)$ karena $!A_g$
merepresentasikan~$g$. Jadi, kita mempunyai
suatu kontradiksi.

Sekarang andaikan $b < \num{m}$. Karena
$\Th{Q} \Proves \lforall[w][(w < \num{m} \lif
\lnot !A_g(w, \num{n}, \Obj 0))]$ menurut \olref{rep-less}, kita memperoleh
$\lnot !A_g(b, \num{n}, \Obj 0)$. Ini kembali bertentangan dengan~(a).
\end{proof}

\end{document}
